\section{Real systems: Gitea, Gogs, GitLab, and a seeded control}
\label{sec:realsystem}

A synthetic system cannot tell us whether the crawl-time reading we identified is
an artefact of our construction, so we repeated the essential comparison on real
systems. We chose Gitea \ResRsVersion{} for genuine object-level sharing with
per-object visibility, a machine-readable sharing surface, and a documented
administrative grant path. Three boundaries on the experiment's conclusions,
stated before the numbers. First, on the endpoints and authorization transitions
exercised in this study, Gitea enforced the expected decisions: no instance of
the studied vulnerability was found, so the experiment measures only the
false-proof side. Second, the cell study uses our instantiation of the
predicates, not the engine; the engine's own relations are executed separately in
Section~\ref{sec:native-run}. Third, the reachability model is reported under
\emph{both} readings---crawl-time snapshot (what the definitions presuppose) and
live at execution time---so the responsible variable can be identified. The
experiment produces \ResRsCells{} cells, \ResRsCellsSnapshot{} per reading.

\paragraph{Grant surface}
Table~\ref{tab:rs-m1} counts an endpoint as grant-capable if, without modifying
the system, a request can move a (subject, object) pair from not-entitled to
entitled. The \ResRsGrantNoAdmin{} endpoints drivable by a non-administrator are
the out-of-band channel in its real form: a subject who is not an administrator
can create an entitlement that a crawl of that subject's own interface need not
reveal. The cross-vendor table below reports a normalised census (semantically
equivalent routes merged, non-authorization grant surface excluded), which is why
its Gitea figure (\GtGrantEndpoints{}) differs from the raw count
(\ResRsGrantEndpoints{}) here; the two answer different questions.

\begin{table}[t]
\centering
\caption{The grant surface of Gitea \ResRsVersion{} (instance-independent).}
\label{tab:rs-m1}
\begin{tabular}{lr}
\toprule
quantity & count \\
\midrule
paths in the system's API description & \ResRsSwaggerPaths{} \\
grant endpoints & \ResRsGrantEndpoints{} \\
revocation endpoints & \ResRsRevokeEndpoints{} \\
\quad of which usable without site-admin & \ResRsGrantNoAdmin{} \\
\bottomrule
\end{tabular}
\end{table}

\paragraph{Is the oracle readable, and from where?}
The oracle's readability depends entirely on the vantage, so the counts below use
different populations and must not be read as one denominator. From the
owner/administrator vantage the sharing surface is readable on
\ResRsOwnerReadable{} of \ResRsOwnerReadableTotal{} attempts: complete. From a
black-box subject's own vantage---the vantage a black-box tool occupies---it is
readable on \ResRsBlackboxSelfReadable{} of \ResRsBlackboxSelfCells{}: never. An
intermediate count of \ResRsSelfReadable{}/\ResRsSelfTotal{} comes only from
subjects who happen to be owners or administrators, and does not constrain the
mechanism; it is therefore not synonymous with ``an arbitrary black-box subject
can read the oracle''. Table~\ref{tab:oracle-readable} collects these
populations. On \ResRsFidelityConsistent{} of \ResRsFidelityTotal{} applicable
cells the credential-substitution check confirmed that our model agrees with the
system's behaviour.

\begin{table}[t]
\centering
\caption{Oracle-readability sampling protocol on Gitea. Each row is a distinct
population and denominator, so the rows are not comparable: both the vantage and
the unit of the denominator differ.}
\label{tab:oracle-readable}
\small
\begin{tabular}{@{}p{0.44\columnwidth}rr@{}}
\toprule
vantage & attempts & readable \\
\midrule
owner / administrator & \ResRsOwnerReadableTotal{} & \ResRsOwnerReadable{} \\
arbitrary subject (black box) & \ResRsBlackboxSelfCells{} & \ResRsBlackboxSelfReadable{} \\
subject also owner/admin (representative grid) & \ResRsSelfTotal{} & \ResRsSelfReadable{} \\
\bottomrule
\end{tabular}
\end{table}

\paragraph{The mismatch is in time, not in visibility}
Table~\ref{tab:rs-m4} is the load-bearing result: the adjudication departs from
the real authorization state on \ResRsMismatchSnapshot{} of \ResRsMismatchDenom{}
cells under the crawl-time reading and \ResRsMismatchLive{} of the same cells
under the live reading---same objects, same entitlements, same predicates. Under
the crawl-time reading, a bookkeeping adjudication produces
\ResRsMrBookkeepingFp{} false proofs on the \ResRsAuthorizedCells{} authorized
cells (rate \ResRsMrBookkeepingFpRate{}), against \ResRsMrBookkeepingAlarmsLive{}
alarms under the live reading; v3 produces \ResRsVThreeFp{} false proofs, and of
the controls \texttt{direct\_only} again over-claims (\ResRsDirectOnlyFp{} false
proofs from \ResRsDirectOnlyAlarms{} alarms) while \texttt{direct\_observe}
emits \ResRsDirectObserveFp{}. Channel provenance of the authorized access
(\ResRsChannelNotObservable{} of \ResRsChannelAuthorizedTotal{} not observable
by the framework log, \ResRsChannelNotObservableRate{}) is a property of our
topology---a design parameter, not a field prevalence.

\begin{table*}[t]
\centering
\caption{Departure from the real authorization state, per reading of the
reachability model.}
\label{tab:rs-m4}
\begin{tabular}{lrr}
\toprule
reading of the reachability model & mismatches & cells \\
\midrule
crawl-time snapshot (as the definitions presuppose) & \ResRsMismatchSnapshot{} & \ResRsMismatchDenom{} \\
live at execution time & \ResRsMismatchLive{} & \ResRsMismatchDenom{} \\
\bottomrule
\end{tabular}
\end{table*}

\paragraph{An endogenous control}
The most serious alternative explanation is that public objects are simply harder
to enumerate---visibility, not time. Two objects that end equally public but
became public at different moments control for this (Table~\ref{tab:rs-endo}):
the predicate reports unreachability on \ResRsEndoAfterCrawlReach{} of
\ResRsEndoAfterCrawlCells{} cells only for the object made public after the
crawl, and on \ResRsEndoFromStartReach{} for the object public throughout, with
the MR firing correspondingly (\ResRsEndoAfterCrawlFires{}
vs.~\ResRsEndoFromStartFires{}). Final visibility is identical; the responsible
variable is when the model was read.

\begin{table*}[t]
\centering
\caption{Endogenous control: two objects that are both public at the end.}
\label{tab:rs-endo}
\begin{tabular}{lrrr}
\toprule
object & cells & \texttt{cannotReach} true & MR fires \\
\midrule
public before the crawl & \ResRsEndoFromStartCells{} & \ResRsEndoFromStartReach{} & \ResRsEndoFromStartFires{} \\
public only after the crawl & \ResRsEndoAfterCrawlCells{} & \ResRsEndoAfterCrawlReach{} & \ResRsEndoAfterCrawlFires{} \\
\bottomrule
\end{tabular}
\end{table*}

\subsection{Cross-vendor replication: Gogs \GgVersion{} and GitLab \GlVersion{}}
\label{sec:gogs}

We repeated the identical protocol on two further systems
(Table~\ref{tab:gogs-cross}): \textbf{Gogs} \GgVersion{}, the project Gitea
forked from in 2016 (related lineage, independently evolved code; census from its
pinned route table, \GgRoutePaths{} paths), and \textbf{GitLab} \GlVersion{}
(no lineage; census from the pinned running instance's routes, \GlRoutePaths{}
paths, since its OpenAPI document omits core grant endpoints). Gogs is a
replication across implementations, GitLab across vendors within one domain.

% 由 analysis/ingest_gitlab.py 自动生成 —— 勿手改
\begin{table}[t]
\centering
\small
\caption{Cross-vendor replication under the identical three-phase protocol
(crawl, out-of-band authorization change, adjudicate) on Gitea
\ResRsVersion{}, Gogs \GgVersion{}, and GitLab \GlVersion{}. Counts are per
reading of the same \ResRsCellsSnapshot{}-cell grid; the live reading is
given in parentheses where both exist. The snapshot reading is the operating
point of the published relations. Gitea and Gogs share lineage
(fork, 2016); GitLab shares none. The team-derived row is read out of the
per-system oracle false negatives by channel (Section~\ref{sec:gogs}). The
Gitea grant count here is the \emph{normalised} M1 census (semantically
equivalent routes merged, non-authorization grant surface excluded), and
therefore differs from the raw route count in Table~\ref{tab:rs-m1}; the two
tables use different normalisations and are not in conflict.}
\label{tab:gogs-cross}
\begin{tabular}{@{}lccc@{}}
\toprule
 & Gitea \ResRsVersion{} & Gogs \GgVersion{} & GitLab \GlVersion{} \\
\midrule
API entries in pinned version (M1 denominator) & 253 & 80 & 1357 \\
Grant endpoints (M1) & 13 & 10 & 18 \\
\;of which usable without site admin & 10 & 4 & 17 \\
Owner-vantage oracle readable & 48 & 48 & 48 \\
\;oracle agrees with actual access & 48 & 31 & 30 \\
\;oracle wrong (false negatives) & 0 & 17 & 18 \\
\;oracle wrong: team-derived access missed & 0 & 4 & 0 \\
Subject-vantage oracle readable & 8 & 35 & 48 \\
Reachability mismatch, snapshot (live) & 10 (0) & 19 (13) & 10 (0) \\
False proofs: bookkeeping alarm, snapshot (live) & 10 (0) & 11 (5) & 10 (0) \\
False proofs: v3 (oracle-gated), snapshot (live) & 0 (0) & 8 (5) & 5 (0) \\
False proofs: direct-only (oracle only) & 0 (0) & 17 (17) & 18 (18) \\
False proofs: direct-observe (oracle $\wedge$ request) & 0 (0) & 17 (17) & 18 (18) \\
\bottomrule
\end{tabular}
\end{table}


The replication changes one thing and confirms two. What it changes concerns the
oracle: \emph{readability does not imply soundness}. Gitea's owner-vantage
surface is readable and sound (agrees with actual access on all \GtOracleAgree{}
cells). Gogs's collaborator-membership endpoint is equally readable
(\GgOracleReadable{}/\GgCells{}) but unsound---wrong on \GgOracleWrong{} cells
(\GgOracleWrongPct\%), exactly the authorizations arising outside direct
membership: site administration (\GgFnSiteAdmin{}), public visibility
(\GgFnPublic{}), out-of-band grants (\GgFnOob{}). GitLab's
\texttt{members/all/\{u\}} resolves inherited membership but public visibility
and site administration confer no membership: \GlOracleWrong{} false negatives
(\GlOracleWrongPct\%). What the oracle cannot see is predictable from what the
surface measures---why clause~2 of the admissibility condition
(Section~\ref{sec:admissibility}) cannot be automated from readability alone.
The consequence is measurable: on Gitea the oracle-gated adjudicators produced no
false proofs (\GtFpVthree{} v3, \GtFpDirectObserve{} direct-observe), while on
Gogs they manufacture them (\GgFpVthree{} snapshot, \GgFpVthreeLive{} live,
\GgFpDirectObserve{} direct-observe) and on GitLab again (\GlFpVthree{},
\GlFpDirectObserve{}), every v3 false proof falling on a visibility-channel
cell. A readable oracle with a blind spot converts its own blind spot into
confirmed violations; on Gogs and GitLab no admissible oracle exists and the
honest adjudicator abstains.

Two results replicate everywhere: bookkeeping alarms false-proof on legitimately
granted access (\GtFpBookkeeping{}/\GgFpBookkeeping{}/\GlFpBookkeeping{} cells)
with the same out-of-band-dominant grant profile; and the reachability mismatch
persists (\GtMismatchSnap{}/\GgMismatchSnap{}/\GlMismatchSnap{} snapshot). On
Gogs the mismatch only partially resolves live (\GgMismatchLive{} remain---team
or public grants never enter a subject's personal crawl); on GitLab the live
reading recovers everything (\GlMismatchLive{} remain)---a second,
vendor-specific failure of the snapshot premise, in the crawl model rather than
the oracle.

\subsection{A positive control with a seeded vulnerability}
\label{sec:seeded}

On the tested endpoints and transitions the three real systems enforce the
expected decisions, so the question remains: \emph{when a real object-level flaw
is present, does the same detector find it?} We answer with a positive control.
A \emph{controlled authorization target} faithfully emulates the Gitea API
contract the engine consumes, with a correct internal authorization model and
\SeedBolaN{} seeded vulnerable instances---all instantiating one fault pattern
(a none-entitlement subject served the object at
\texttt{GET /api/v1/repos/\{owner\}/\{repo\}}), differing in subject and object,
not in fault location; Table~\ref{tab:seeded} is therefore evidence about this
fault pattern, not three independent defect discoveries. Two design points: the
target deliberately exposes its authorization state to the subject-side harness
(mirroring the admissibility condition), which production Gitea does not---on the
seeded target the subject-side harness can read the authorization state of every
cell by design, whereas on production Gitea the oracle is readable from the
owner/admin vantage (\ResRsOwnerReadable{}/\ResRsOwnerReadableTotal{}) but from an
arbitrary subject's vantage it is not (\ResRsBlackboxSelfReadable{}/\ResRsBlackboxSelfCells{});
and the introspection surface
reports the \emph{correct} \texttt{none} permission, so the oracle stays sound
and the defect is visible only as a discrepancy between oracle and request. We
drove it with the same topology, mutation, crawl, introspection and adjudication
code as above, differing only in the system under test; ground truth comes from
the seed list, not from the detector.

\begin{table}[t]
\centering
\caption{Positive control: detection of \SeedBolaN{} seeded vulnerable instances
(one fault pattern). TP/FN/FP/TN are counted against the known seed, not against
a model.}
\label{tab:seeded}
\begin{tabular}{lrrrr}
\toprule
method & TP & FN & FP & TN \\
\midrule
\texttt{mr\_bookkeeping} & \SeedTPMrBookkeeping{} & \SeedFNMrBookkeeping{} & \SeedFPMrBookkeeping{} & \SeedTNMrBookkeeping{} \\
\texttt{v3} (owner vantage) & \SeedTPVThree{} & \SeedFNVThree{} & \SeedFPVThree{} & \SeedTNVThree{} \\
\texttt{v3} (subject vantage) & \SeedTPVThreeSelf{} & \SeedFNVThreeSelf{} & \SeedFPVThreeSelf{} & \SeedTNVThreeSelf{} \\
\texttt{direct\_only} & \SeedTPDirectOnly{} & \SeedFNDirectOnly{} & \SeedFPDirectOnly{} & \SeedTNDirectOnly{} \\
\texttt{direct\_observe} & \SeedTPDirectObserve{} & \SeedFNDirectObserve{} & \SeedFPDirectObserve{} & \SeedTNDirectObserve{} \\
\texttt{direct\_observe} (subject) & \SeedTPDirectObserveSelf{} & \SeedFNDirectObserveSelf{} & \SeedFPDirectObserveSelf{} & \SeedTNDirectObserveSelf{} \\
\bottomrule
\end{tabular}
\end{table}

The observation-gated adjudicator behaves exactly as the negative result predicts
(Table~\ref{tab:seeded}): \texttt{direct\_observe} and \texttt{v3} catch all
\SeedBolaN{} planted instances with zero false positives; \texttt{direct\_only}
again over-claims (\SeedFPDirectOnly{} false positives on cells where the oracle
says \texttt{none} but no access succeeds). The control narrows the claim to this
fault pattern: when a sound oracle is available, the metamorphic step is not
necessary for detection here---both \texttt{direct\_observe} and \texttt{v3}
catch all \SeedBolaN{} instances, while \texttt{direct\_only} again over-claims.
It does not generalise beyond one fault pattern, and the seeded BFLA defect
(below) remains outside the read-oriented logic's coverage.
The control also delimits the detector: we seeded one BFLA defect (non-privileged
\SeedBflaActor{} succeeding at \texttt{PUT \dots/collaborators/\{u\}}, HTTP
\SeedBflaStatus{}), which the read-oriented logic does not cover---a known
boundary, reported as such rather than as a detection failure. This is a
controlled research artefact, not production Gitea; it validates the detector's
logic, not Gitea's CVE posture. Its source, seed list and run script are in the
reproducibility package at \url{https://github.com/ubfbuz3/JISA} (tag
\texttt{artifact-v1}).
